\documentclass[UTF8,a4paper,10pt]{ctexart}

% 保持课程 Overleaf 模板的版心、字体层级和页眉页脚结构。
\usepackage[left=3.17cm,right=3.17cm,top=2.74cm,bottom=2.74cm,headheight=15pt]{geometry}
\usepackage{amsmath,amssymb}
\usepackage{graphicx,subfig,float,cite,caption,enumerate}
\usepackage{booktabs,longtable,tabularx,array,multirow}
\usepackage{fancyhdr}
\usepackage{algorithm,algorithmicx,algpseudocode}
\usepackage{xcolor,listings,hyperref,enumitem,pifont,tikz,eso-pic}
\usepackage{fontspec}

\definecolor{nkublue}{HTML}{244A73}
\definecolor{lightblue}{HTML}{EAF0F6}
\definecolor{codebg}{HTML}{F5F7F9}
\definecolor{codegray}{HTML}{4B5563}
\hypersetup{
  colorlinks=true,
  linkcolor=nkublue,
  citecolor=nkublue,
  urlcolor=nkublue,
  pdftitle={预备实验：编译流程观察与 RISC-V 汇编编程},
  pdfauthor={马禹翔；胡家华}
}

\IfFontExistsTF{Noto Serif CJK SC}{
  \setCJKmainfont[ItalicFont=Noto Sans CJK SC Bold,BoldFont=Noto Serif CJK SC Black]{Noto Serif CJK SC}
}{
  \setCJKmainfont[BoldFont=SimHei]{SimSun}
}

\ctexset{
  section={name={,、},number={\chinese{section}}},
  subsection={name={（,）},number={\chinese{subsection}}},
  subsubsection={name={,.},number={\arabic{subsubsection}}}
}

\setlength{\parindent}{2em}
\setlength{\parskip}{0.15em}
\setlength{\emergencystretch}{3em}
\setlist{nosep,leftmargin=2em}
\renewcommand{\arraystretch}{1.25}
\captionsetup{font=small,labelsep=quad}

\pagestyle{fancy}
\fancyhf{}
\lhead{\kaishu \leftmark}
\rhead{\kaishu 编译原理实验报告}
\fancyfoot[C]{\thepage}
\renewcommand{\headrulewidth}{0.4pt}
\renewcommand{\footrulewidth}{0pt}

\lstdefinestyle{code}{
  basicstyle=\ttfamily\small,
  backgroundcolor=\color{codebg},
  frame=single,
  rulecolor=\color{lightblue},
  numbers=left,
  numberstyle=\scriptsize\color{codegray},
  numbersep=8pt,
  breaklines=true,
  columns=fullflexible,
  keepspaces=true,
  showstringspaces=false,
  tabsize=4,
  xleftmargin=1.5em,
  framexleftmargin=1em
}
\lstset{style=code}

\newcolumntype{Y}{>{\raggedright\arraybackslash}X}

\newcommand{\HRule}{\rule{\linewidth}{0.5mm}}
\newcommand\watermark[3]{%
  \AddToShipoutPictureBG{%
    \parbox[b][\paperheight]{\paperwidth}{%
      \vfill\centering
      \begin{tikzpicture}
        \node[rotate=#1] {\includegraphics[width=0.18\paperwidth]{#3}};
      \end{tikzpicture}
      \vfill
    }%
  }%
}

\begin{document}
\begin{titlepage}
\begin{center}
\includegraphics[width=0.8\textwidth]{NKU.png}\\[1cm]
\textsc{\Huge \kaishu{\textbf{南 开 大 学}} }\\[0.9cm]
\textsc{\huge \kaishu{\textbf{计\ \ 算\ \ 机\ \ 学\ \ 院}}}\\[0.5cm]
\textsc{\Large \textbf{编译原理实验报告}}\\[0.8cm]
\HRule \\[0.9cm]
{\LARGE \bfseries 了解编译器：LLVM IR 编程与 RISC-V 汇编编程}\\[0.4cm]
\HRule \\[1.7cm]
\centering
\textsc{\LARGE \kaishu{马禹翔 2410933}}\\[0.5cm]
\textsc{\LARGE \kaishu{胡家华 2410935}}\\[0.5cm]
\textsc{\LARGE \kaishu{指导教师：王刚}}\\[0.5cm]
\vfill
{\Large 2026 年 10 月 10 日}
\end{center}
\end{titlepage}

\clearpage
\pagenumbering{Roman}

\begin{abstract}
本实验围绕编译器完整处理流程，将同一 SysY 程序分别表示为 LLVM IR 和 RISC-V 汇编，并链接课程提供的 SysY 运行库执行。实验以斐波那契程序对齐小组两人的成果，另设计包含整数运算、短路逻辑、作用域、循环控制和函数调用的补充程序，并扩展一维数组与单精度运算。采用 Clang/LLVM 生成 RV64GC、LP64D 静态 ELF，在 QEMU 中执行，以源程序编译结果和独立 Python 算法共同核验。共完成 110 组输入、570 次执行，标准输出及退出状态全部符合预期。实验还比较了 \texttt{-O0}、\texttt{-O2} 与调试选项产生的中间文件，说明目标设置、调用约定和运行库兼容性对跨架构整合的影响。

\end{abstract}

\clearpage
\tableofcontents
\clearpage
\pagenumbering{arabic}
\watermark{60}{10}{NKU-watermark.png}

\begin{table}[H]
  \centering
  \caption*{小组分工}
  \begin{tabularx}{\textwidth}{>{\bfseries}p{3.0cm}Y}
    \toprule
    成员 & 分工及本报告采用的成果 \\
    \midrule
    马禹翔（2410933） & RISC-V 汇编分工：斐波那契汇编、语言特性补充例、运行库链接、交叉验证与报告整合。 \\
    胡家华（2410935） & LLVM 分工：原始斐波那契 C/SysY、手写 LLVM IR、LLVM 观察记录及 MLIR 文档阅读探索。 \\
    \bottomrule
  \end{tabularx}
\end{table}

合作仓库基准为 \href{https://github.com/hjhuiuaa/Compiler}{hjhuiuaa/Compiler} 的“预备工作”分支.
\section{引言}

高级语言中的变量、表达式和控制流，需要逐步转化为机器指令与可执行文件。本实验的重点是验证这些表示之间的对应关系，而不止于生成一份汇编文本。我们将任务分成两条路径：一条保存预处理、抽象语法树、LLVM IR、目标汇编和 ELF 的阶段证据；另一条验证手写 LLVM IR 与手写 RISC-V 汇编的可观察行为。两条路径最终使用同一目标架构、同一运行库和相同输入汇合。

课程基础要求涉及整数类型、常量与变量、赋值、条件分支、循环、函数、作用域、\texttt{break}/\texttt{continue}、逻辑运算及 SysY 运行库连接。指导书给出的斐波那契程序适合观察循环，却不能单独覆盖所有特性。因此，本实验保留斐波那契作为小组共同语义基准，同时加入语言特性补充例.

\section{实验环境与编译处理流程}

\subsection{实验环境}

\begin{table}[H]
  \centering
  \caption{整合后重新验证的实验环境}\label{tab:env}
  \begin{tabularx}{\textwidth}{p{3.4cm}Y}
    \toprule
    项目 & 实际配置 \\
    \midrule
    宿主与工具环境 & Windows；独立 WSL2 Alpine Linux 3.22 \\
    编译器与链接器 & Alpine Clang 20.1.8；LLD 20.1.8 \\
    执行器 & qemu-riscv64 10.0.0 \\
    目标与 ABI & \texttt{riscv64-linux-musl}；\texttt{-march=rv64gc}；\texttt{-mabi=lp64d} \\
    标准库与 SysY 库 & RISC-V musl 1.2.5-r12；课程 \texttt{sylib.c} 重新编译得到的静态库 \\
    \bottomrule
  \end{tabularx}
\end{table}

\subsection{语言处理系统的各阶段}

Clang 命令本身是驱动程序，负责调度前端、优化器、目标代码生成器、汇编器和链接器。不同停止选项可以保存各阶段输出\cite{clang-guide}。

\begin{longtable}{p{2.6cm}p{3.9cm}p{6.7cm}}
  \caption{编译阶段、输出及本次观察内容}\label{tab:stages}\\
  \toprule
  阶段 & 保存的输出 & 主要观察 \\
  \midrule
  \endfirsthead
  \toprule
  阶段 & 保存的输出 & 主要观察 \\
  \midrule
  \endhead
  预处理 & \texttt{fib.i}、\texttt{fib.tokens.txt} & 展开头文件和宏，处理条件编译与行标记，得到记号序列。 \\
  语法与语义分析 & \texttt{fib.ast.txt} & 观察 \texttt{VarDecl}、\texttt{WhileStmt}、\texttt{CallExpr} 及表达式类型。 \\
  IR 生成与优化 & \texttt{fib.O0.ll}、\texttt{fib.O2.ll} & 观察栈槽、\texttt{load/store}、比较、分支和 SSA 合并。 \\
  目标代码生成 & \texttt{fib.O0.s}、\texttt{fib.O2.s} & 观察指令选择、寄存器分配、栈帧与函数调用。 \\
  汇编 & \texttt{*.o}、\texttt{object.txt} & 将助记符编码为指令，建立节、符号与待处理重定位。 \\
  链接 & ELF、\texttt{*.map} & 合并目标文件与库，解析外部符号并布局最终地址。 \\
  \bottomrule
\end{longtable}

核心命令如下；其中 \texttt{\$TARGET} 表示目标三元组、ABI 和 sysroot 参数，完整命令及退出码保存在 \texttt{build/commands.json}。

\begin{lstlisting}[language=bash,caption={编译流程中的阶段命令}]
clang $TARGET -E fib.c -o fib.i
clang $TARGET -fsyntax-only -Xclang -ast-dump fib.c
clang $TARGET -O0 -S -emit-llvm fib.c -o fib.ll
clang $TARGET -O0 -S fib.c -o fib.s
clang $TARGET -c fib.s -o fib.o
\end{lstlisting}

观察 C 程序时使用 \texttt{stdio} 接口；手写 IR 和汇编使用 \texttt{getint}、\texttt{putint}、\texttt{putch}。两者的算法和标准输出相同。

\section{示例程序与语言特性覆盖}

\subsection{斐波那契：小组共同的语义基准}

沿用合作仓库中的 \texttt{fib.sy}：先输出 0 和 1，令 $i=1$；当 $i<n$ 时执行
$t=b$、$b=a+b$、输出 $b$、$a=t$、$i=i+1$。输入 5 时得到 0、1、1、2、3、5，共六行。$n\leq 1$ 时仍然输出前两项，这是指导书程序本身的行为，因此没有擅自改成“输出 $n$ 项”。测试限制到 $n\leq 46$，因为 $F(46)=1836311903$ 仍在有符号 32 位整数范围内，而 $F(47)$ 会溢出。

\subsection{补充例的设计}

\texttt{examples/features.sy} 读取整数 $x$、$d$，测试约定二者均位于 $[-100,100]$。程序依次输出加、减、乘、除、余、符号分类、逻辑结果、副作用计数、数组归约、浮点转整数和作用域观察值。$d=0$ 时显式避开除法与取模。

\begin{table}[H]
  \centering
  \caption{补充例对语言特性的覆盖}
  \begin{tabularx}{\textwidth}{p{4.0cm}Y}
    \toprule
    特性 & 示例位置或验证方式 \\
    \midrule
    常量、变量、初始化、赋值 & \texttt{WIDTH}、\texttt{calls}，局部 $x/d$，\texttt{reduce} 中的 $i/sum$ 更新 \\
    算术与单目运算 & 主函数前五项结果及 \texttt{emit(-x)} \\
    六种关系运算 & 符号分类、循环条件、短路条件与逻辑非组合 \\
    逻辑与、逻辑或、逻辑非 & 除数保护、\texttt{touch} 副作用与 \texttt{!(x <= d)} \\
    条件、循环和控制转移 & \texttt{if/else}、\texttt{while}、\texttt{break}、\texttt{continue} \\
    函数与返回 & \texttt{touch}、\texttt{reduce}、\texttt{blend}、\texttt{emit} \\
    嵌套作用域 & 内层 \texttt{int x=d+1}，离开语句块后再次输出外层 $x$ \\
    运行库输入输出 & \texttt{getint} 输入，\texttt{putint/putch} 输出 \\
    额外覆盖 & 一维整数数组与数组参数；\texttt{float} 运算及隐式转换 \\
    \bottomrule
  \end{tabularx}
\end{table}

数组初始化为 \texttt{\{x,-2,d,0,9\}}。\texttt{reduce} 先检查零，再跳过负数，因此尾部 9 永远不累加，可以同时检验 \texttt{break} 和 \texttt{continue}。浮点表达式为
\[
z=\frac{(x+0.5)\times 2-d}{2},
\]
转成整数时向零截断。测试范围内的中间半整数都可以精确表示，因此这里不需要浮点误差容限。

SysY 是 C 的子集，构建脚本用 \texttt{-x c} 和单独的接口声明编译这些 \texttt{.sy} 文件作为参考实现；这不代表 Clang 内建了 SysY 前端。本实验也没有声称覆盖多维数组、张量扩展或实现完整 SysY 编译器。

\section{LLVM IR 成果整合}

\subsection{原始成果}

\texttt{fib.ll} 是 Clang 15 自动生成的参考 IR，\texttt{fib-sysy.ll} 是将输入输出替换成 SysY 运行库后的手写版本。原始 IR 将 $a,b,i,t,n$ 放在 \texttt{i32} 栈槽中；赋值对应 \texttt{store}，读取变量对应 \texttt{load}；循环条件使用 \texttt{icmp slt}，随后由 \texttt{br} 选择循环体或出口。输入使用 \texttt{call i32 @getint()}，每个整数的输出由 \texttt{putint} 与 \texttt{putch(10)} 组成。

记录到数字 SSA 名称不连续导致 \texttt{llvm-as} 报错的问题。数字结果名与数字基本块标签共享编号顺序，\texttt{void} 调用不占结果编号；移除原先的 \texttt{printf} 调用以后，需要同步调整后续数字名称。补充 IR 改用语义化名称，以降低维护负担。LLVM 的 \texttt{ptr} 为不透明指针，访问数据的类型由 \texttt{load/store} 等指令指定\cite{llvm-langref}。

\subsection{由 x86 目标移植到 RISC-V}

原始文件带有 \path{x86_64-pc-linux-gnu} 目标三元组、x86 数据布局和
\path{target-cpu/target-features/tune-cpu} 属性，不能只更改文件扩展名就当作 RISC-V 输入。本次保留原文件，由 \path{scripts/port_ir.py} 生成 \path{ir/fib-riscv.ll}：替换目标三元组，移除旧数据布局与 x86 专属属性，并删除不再调用的 \texttt{stdio} 声明。Clang 编译时再以显式目标参数确定实际 Linux musl 数据布局。

斐波那契的基本块结构和 \texttt{i32} 语义没有改变。\texttt{llvm-as} 将文本 IR 转成 bitcode，\texttt{opt -passes=verify} 检查结构，之后经 Clang 后端生成 RV64 目标文件。该程序没有指针与整数互转等依赖原数据布局的操作，适合这样移植；这一经验不能无条件推广到任意 LLVM IR。

\subsection{补充程序的 IR 设计}

\texttt{ir/features-riscv.ll} 是新增的手写实现。\texttt{reduce} 通过 \texttt{phi} 表示循环携带的 $i$ 与 $sum$，\texttt{getelementptr} 按 \texttt{i32} 元素寻址；主函数将短路求值拆成多个基本块，确保 $d=0$ 时不会执行 $x/d$；浮点部分使用 \texttt{fadd/fmul/fsub/fdiv}，再由 \texttt{fptosi} 转换为整数。

\begin{lstlisting}[caption={逻辑与的短路控制流}]
and.left:
  br i1 %nonzero, label %and.right, label %and.false
and.right:
  %q2 = sdiv i32 %x, %d
  %gt = icmp sgt i32 %q2, 1
  br i1 %gt, label %and.true, label %and.false
\end{lstlisting}

不能提前计算带除法的右操作数再执行普通按位与：即使最终布尔值为假，提前求值仍然会改变源程序行为。同理，逻辑或右侧的 \texttt{touch} 只能在左侧为假时调用。

\section{RISC-V 汇编实现}

\subsection{调用约定与数据宽度}

本实验采用 RV64GC/LP64D。整数参数和返回值通过 \texttt{a0--a7} 传递，\texttt{blend} 的两个浮点参数使用 \texttt{fa0}、\texttt{fa1}。跨函数调用后仍需存活的整数放在 \texttt{s} 类寄存器中，函数在入口保存、出口恢复；\texttt{ra} 也必须保存，因为 \texttt{call} 会覆盖它。所有调用处的 \texttt{sp} 均保持 16 字节对齐，符合 RISC-V psABI\cite{riscv-abi}。

\begin{table}[H]
  \centering
  \caption{手写汇编的寄存器与栈帧安排}
  \begin{tabularx}{\textwidth}{p{3.7cm}Y}
    \toprule
    程序或函数 & 存储分配 \\
    \midrule
    \texttt{fib} 的 \texttt{main} & \texttt{s0=a}、\texttt{s1=b}、\texttt{s2=i}、\texttt{s3=n}；48 字节栈帧保存 \texttt{ra} 和 \texttt{s0--s3} \\
    \texttt{features} 的 \texttt{main} & \texttt{s0=x}、\texttt{s1=d}；64 字节栈帧；\texttt{sp+0} 起存放五个 \texttt{i32} 元素 \\
    \texttt{emit} & 16 字节栈帧保存 \texttt{ra}，依次调用 \texttt{putint} 和 \texttt{putch} \\
    reduce、touch、blend & 叶函数，使用临时或参数寄存器，无需建立栈帧 \\
    \bottomrule
  \end{tabularx}
\end{table}

RV64 的整数寄存器宽度是 64 位，但 SysY \texttt{int} 仍为 32 位。因此整数计算使用
\texttt{addw/subw/mulw/divw/remw}，数组访存使用 \texttt{lw/sw}；保存 \texttt{ra} 和 \texttt{s} 寄存器则使用 \texttt{ld/sd}。

\subsection{斐波那契循环}

\begin{lstlisting}[caption={手写斐波那契循环}]
.Lfib_cond:
    bge s2, s3, .Lfib_exit
    mv t0, s1
    addw s1, s0, s1
    mv s0, t0
    mv a0, s1
    call putint
    li a0, 10
    call putch
    addiw s2, s2, 1
    j .Lfib_cond
\end{lstlisting}

循环头保持不变式 \texttt{s0=F(i-1)}、\texttt{s1=F(i)}。一次更新后输出 $F(i+1)$，再递增 $i$，所以执行顺序与 SysY 源程序一致。源程序中的临时变量 $t$ 放在 \texttt{t0}；在下一次函数调用前，它已转移到保存寄存器 \texttt{s0}，因而不会错误依赖调用者保存的临时寄存器。

\subsection{补充汇编的关键细节}

逻辑与首先用 \texttt{beqz} 检查 $d$，逻辑或在 $d=0$ 时直接跳过 \texttt{touch}。数组元素地址由基址加 $i\times4$ 得到。浮点部分使用 \texttt{flw}、\texttt{fadd.s}、\texttt{fmul.s}、\texttt{fsub.s}、\texttt{fdiv.s}；\texttt{fcvt.w.s} 明确选择 \texttt{rtz}（向零舍入）。若使用默认舍入模式，半整数用例会偏离语言规定的浮点转整数行为。内层同名变量 $x$ 仅暂存在参数寄存器中，外层 $x$ 始终保留于 \texttt{s0}，实现作用域隔离。

\section{汇编、链接与运行验证}

\subsection{链接课程 SysY 运行库}

本实验使用目标 Clang 重新编译 \texttt{sylib.c}，再以 \texttt{llvm-ar} 建立 \path{build/libsysy_riscv.a}。程序与库使用相同的目标架构、ABI 与 RISC-V musl 标准库；没有链接宿主 x86 库，也没有用自制输入输出函数代替课程运行库。

实验采用静态链接，显式加入 musl 启动文件、\texttt{libc.a} 与目标 \texttt{compiler-rt} 内建函数库。首次链接出现 \texttt{\_\_addtf3} 等未定义符号，原因是 libc 格式化输入输出需要四倍精度辅助运算；补入 RISC-V builtins 后解决，初次诊断保存在 \path{evidence/initial-link-missing-builtins.json}。

\texttt{fib-asm.o} 中的 \texttt{getint/putint/putch} 是未定义符号，调用位置带 RISC-V 重定位；链接完成后，最终符号表中出现它们的定义。\texttt{readelf} 确认最终程序为 RISC-V ELF64 静态可执行文件，ELF 标志包含 RVC 和 double-float ABI，QEMU 可以直接运行。重定位、反汇编、符号表和链接映射都保存在 \texttt{build/stages}。

\subsection{测试方法与结果}

\begin{table}[H]
  \centering
  \caption{差分测试与独立参考测试结果}
  \begin{tabularx}{\textwidth}{p{2.6cm}Yp{3.0cm}}
    \toprule
    测试集 & 输入与版本 & 实测结果 \\
    \midrule
    斐波那契 & $-5,0,1,2,3,5,10,20,30,46$；7 个版本 & $10\times7=70$ 次通过 \\
    语言特性 & 49 组符号/零/边界组合和 51 组固定种子随机输入；5 个版本 & $100\times5=500$ 次通过 \\
    汇总 & 110 组输入；570 次独立进程执行 & 570 通过，0 失败 \\
    \bottomrule
  \end{tabularx}
\end{table}

五个基本版本分别为 SysY 源程序的 \texttt{-O0/-O2}、手写 IR 的 \texttt{-O0/-O2} 与手写汇编；斐波那契还执行原始 \texttt{fib.c} 的 \texttt{-O0/-O2}。每个版本都与独立 Python 算法生成的预期字符串逐字比较，并要求退出码为 0。

输入 \texttt{7 3} 时，补充程序依次输出
\[
10,4,21,2,1,1,1,1,1,1,10,6,-7,4,7.
\]
输入 \texttt{7 0} 时，副作用计数为 0，说明逻辑或确实跳过 \texttt{touch}；浮点结果为 7，检验了 7.5 向零截断。负数测试用于检查有符号比较、除法和余数语义，测试预言没有直接使用 Python 的 \texttt{//} 与 \texttt{\%}，因为它们对负数的规则与 C/SysY 不同。

测试将标准输出和标准错误分开保存。运行库的 \texttt{TOTAL: ...} 行属于标准错误，不是程序数值输出，也没有被用于性能评价。完整输入、预期值、实际输出和退出码见 \texttt{build/test-results.json}。有限测试证明这些用例上的一致性，但不等价于对所有输入的形式化证明。

\section{优化选项与结果讨论}

\subsection{优化和调试选项的影响}

\begin{table}[H]
  \centering
  \caption{\texttt{fib.c} 在不同优化级别下的真实输出统计}
  \begin{tabular}{lrr}
    \toprule
    观察指标 & \texttt{-O0} & \texttt{-O2} \\
    \midrule
    IR 中 \texttt{alloca} 数量 & 6 & 1 \\
    IR 中 \texttt{load} 数量 & 10 & 2 \\
    IR 中 \texttt{store} 数量 & 8 & 0 \\
    IR 中 \texttt{phi} 数量 & 0 & 3 \\
    目标文件 \texttt{text} 列字节数 & 247 & 197 \\
    \bottomrule
  \end{tabular}
\end{table}

\texttt{-O0} 便于将源变量对应到栈槽；\texttt{-O2} 将可提升的变量合并为 SSA 值，减少内存访问，并改变控制流及寄存器保存安排。\texttt{scanf} 需要取得 $n$ 的地址，因此 $n$ 的存储需求与其他局部变量不同。表中的 \texttt{text} 是 \texttt{llvm-size} 的统计口径，包含只读内容等，不能解释成 \texttt{main} 单个函数的纯指令长度；这里只做同一源码配置间的比较，不据此推导运行速度。

\texttt{fib} IR 保留了 \texttt{optnone}，所以给编译命令加入 \texttt{-O2} 不能解释成“对该函数完成了完整中端优化”。这两个版本用于检查不同编译选项下仍能生成可执行文件；优化现象比较使用的是重新从 \texttt{fib.c} 生成的 IR。加入 \texttt{-g} 后，目标文件出现 \texttt{.debug\_info}、\texttt{.debug\_line} 等调试节，便于将机器地址映射回源位置；调试信息和程序正确性属于不同维度。

\subsection{跨架构整合经验}

LLVM IR 并非天然完全不依赖目标：目标三元组、数据布局、函数属性和调用约定必须一致。汇编文本也不等于可执行文件：外部函数地址由链接阶段确定，运行库还依赖 C 标准库和启动代码。本实验生成面向 Linux 应用程序的静态 ELF，没有机械照搬讲义中针对另一套工具链的固定链接地址；\texttt{-mcmodel=medany} 和 \texttt{--no-relax} 均明确记录在构建参数中。

最终的 58 条构建与检查命令都返回 0。日志中 \texttt{overriding the module target triple} 表示构建参数将 IR 中的 Linux GNU 标记明确覆盖为 Linux musl；对已经生成的汇编文件传入共同编译参数时还会产生少量“参数未使用”警告。这些诊断全部保留在日志中，但不代表构建失败。

\section{进阶要求：MLIR 与 AscendNPU IR}

\subsection{多级 IR 与语义逐层下降}

LLVM IR 已经采用 SSA、基本块和接近机器的数据操作，适合通用优化及后端代码生成；但卷积、矩阵乘、张量形状等高层语义在下降为 \texttt{load/add/store} 后难以恢复。NPU 还具有 Cube/Vector 两类计算核以及 GM、UB、L1 等存储层次，通用 LLVM IR 中没有直接对应的抽象。课程提供的 AscendNPU IR 资料据此说明 MLIR 的设计动机：不同抽象层使用不同方言表达，通过 Lowering 在保持程序语义的前提下逐步引入硬件信息。

MLIR 的基本构成包括 Dialect、Operation 和 Type。Dialect 提供操作与类型的命名空间；Operation 是 IR 的基本节点，可以携带操作数、结果和编译期属性；Type 描述值或内存对象的类型。不同方言可以同时存在于一个模块中，因此一次 Lowering 不必把整个程序直接改写成最终指令，而可以只替换当前阶段已经合法化的操作。

\begin{table}[H]
  \centering
  \caption{AscendNPU IR 三个核心方言的职责划分}
  \begin{tabularx}{\textwidth}{p{2.1cm}p{2.2cm}Yp{3.2cm}}
    \toprule
    方言 & 抽象层次 & 主要表达内容 & 代表性操作或属性 \\
    \midrule
    HFusion & 高层、硬件无关 & 张量级算子语义，以及算子融合、形状传播和自动调度 & \texttt{hfusion.vadd}、\texttt{hfusion.matmul} \\
    HIVM & NPU 硬件抽象 & NPU 核间同步、片上内存分配、流水向量化，以及 GM/UB 地址空间 & \texttt{hivm.hir.load}、\texttt{hivm.hir.vadd}、\texttt{hivm.hir.store} \\
    HACC & 异构执行抽象 & Host/Device 启动语义、设备入口和 tiling 函数绑定 & \texttt{hacc.entry}、设备函数属性 \\
    \bottomrule
  \end{tabularx}
\end{table}

这三层并不是互相替代的三种写法。HFusion 保留便于融合和调度的算子语义；HIVM 将算子映射到 NPU 存储与计算资源；HACC 说明设备函数如何进入异构运行时。此后才下降到 LLVM IR intrinsic，由后端完成寄存器分配、指令调度与机器码生成\cite{ascend-arch}。

\subsection{VecAdd 的逐层 Lowering 分析}

课件以 16 个 \texttt{i16} 元素的向量加法为贯穿示例。前端 Arith/Tensor IR 只表达“两个张量相加”，尚未确定内存位置。\texttt{arith2fusion} 或相应转换首先把值语义的 \texttt{tensor} 转成内存语义的 \texttt{memref}，并用 \texttt{hfusion.vadd} 保存高层向量加法语义。随后 \texttt{convert-to-hivm-op} 将一个高层操作展开为数据搬入、计算和结果写回，同时为中间缓冲标注 UB 地址空间。

\begin{lstlisting}[caption={VecAdd 从 HFusion 到 HIVM 的结构变化}]
// HFusion: one tensor-level operation
%c = hfusion.vadd ins(%a, %b : memref<16xi16>)
                     outs(%cbuf : memref<16xi16>)

// HIVM: explicit storage and data movement
%ub_a = memref.alloc() : memref<16xi16, #hivm.address_space<ub>>
%ub_b = memref.alloc() : memref<16xi16, #hivm.address_space<ub>>
%ub_c = memref.alloc() : memref<16xi16, #hivm.address_space<ub>>
hivm.hir.load  ins(%a_gm)  outs(%ub_a)
hivm.hir.load  ins(%b_gm)  outs(%ub_b)
hivm.hir.vadd  ins(%ub_a, %ub_b) outs(%ub_c)
hivm.hir.store ins(%ub_c) outs(%c_gm)
\end{lstlisting}

Lowering 后出现的 \texttt{gm} 与 \texttt{ub} 不是注释，而是后续内存规划和指令选择需要的类型信息。对于该示例，每个向量占 32 字节，三个 UB 缓冲区可依次安排在偏移 0、32 和 64。\texttt{Plan-Memory} 将抽象的 \texttt{memref.alloc} 具体化为不重叠的地址区间；\texttt{Inject-Sync} 在数据搬运与向量计算之间插入等待，避免 \texttt{vadd} 在 \texttt{load} 完成前读取数据。更大的张量还可以通过 Multi-Buffer 让当前分块计算与下一分块搬运重叠，从而隐藏访存延迟。

\begin{table}[H]
  \centering
  \caption{VecAdd 下降过程中的信息变化}
  \begin{tabularx}{\textwidth}{p{2.2cm}p{3.2cm}Y}
    \toprule
    层级 & 主要表示 & 新增或保留的信息 \\
    \midrule
    Level 0 & Arith/Tensor IR & 张量类型和逐元素加法，不包含设备信息 \\
    Level 1 & HFusion IR & NPU 编译流程可识别的高层算子，便于合法化、化简和融合 \\
    Level 2 & HIVM IR & GM/UB 地址空间、显式搬运、片上缓冲和同步关系 \\
    Level 3 & HACC 语义 & 设备入口、Host/Device 函数类别和启动关系 \\
    Level 4 & LLVM IR & \texttt{llvm.npu.*} intrinsic，交由 LLVM 后端生成设备二进制 \\
    \bottomrule
  \end{tabularx}
\end{table}

这一过程说明 Lowering 的关键是把隐含约束转成显式 IR。高层的一次加法在 HIVM 中成为两次 GM 到 UB 的搬运、一次向量运算和一次 UB 到 GM 的写回；到了 HACC 层，函数还要带有设备入口属性。每一层仍表达同一计算，但能够执行的分析和优化已经不同。

\subsection{Pass 流水线、工具链与验证范围}

课件给出的完整思路可以概括为：前端方言先经过 Legalize 与 Decompose；HFusion 阶段执行算子化简、融合和自动调度；转入 HIVM 后进行内存规划、同步插入、Mix Kernel 拆分与多缓冲；HACC 绑定异构执行语义，最后由 \texttt{hivmc} 转成 LLVM IR 并生成 \texttt{kernel.o}。其中 \texttt{Split-Mix-Kernel} 对同时使用 Cube 与 Vector 的算子有意义；纯向量 VecAdd 只需映射到 Vector 核，不发生双核拆分。

\texttt{bishengir-opt} 可单独运行 Pass，并利用 MLIR 的打印选项观察每一步变化。若具备 CANN 与毕昇编译器环境，可以采用以下命令保存中间 IR：

\begin{lstlisting}[language=bash,caption={逐 Pass 观察 AscendNPU IR 的命令形式}]
bishengir-opt add.mlir \
  --convert-hfusion-to-hivm \
  --hivm-plan-memory \
  --hivm-inject-sync \
  --mlir-print-ir-after-all 2>&1 | tee after-each-pass.mlir

bishengir-translate --mlir-to-llvmir add.mlir > add.ll
bishengir-compile add.mlir -enable-hivm-compile -o kernel.o
\end{lstlisting}

本实验对上述内容完成了课件、官方架构说明和 VecAdd 示例之间的对应分析\cite{ascend-example,ascend-arch}。当前实验环境没有安装 CANN，也没有可用的昇腾 NPU 或 Atlas 开发板，因此没有把这些命令的预期输出写成实测结果。基础部分的 570 次执行只统计 RISC-V ELF 在 QEMU 中的真实运行；进阶部分给出可复现的工具链路径和逐层 IR 判读依据，二者的证据范围分别说明。

\section{结论与复现方法}

本实验完成了预处理、前端分析、LLVM IR、目标汇编、目标文件及链接后 ELF 的证据链。手写 RISC-V 程序通过课程运行库完成输入输出，与移植后的手写 LLVM IR 以及源程序编译结果在 110 组输入上保持一致。汇编部分同时验证了 32 位整数语义、64 位寄存器环境、调用保存规则、短路副作用和浮点向零截断等容易出错的边界。

\begin{lstlisting}[language=bash,caption={实验复现命令}]
python3 scripts/build.py
python3 scripts/test.py
echo 5 | qemu-riscv64 build/fib-asm
printf '7 3\n' | qemu-riscv64 build/features-asm
\end{lstlisting}

\bibliographystyle{plain}
\bibliography{references}

\end{document}
